\chapter{구면, 쌍곡 공간, 그리고 \(k\)의 의미}

간단한 수학을 조금 해 봅시다. 차원이 몇이든 구는 특별한 기하학적 대상입니다. 구는 균일하며, 여기서 균일하다는 것은 곡률이 일정하다는 뜻입니다. 먼저 양의 곡률을 살펴보고, 이어서 균일한 음의 곡률을 본 다음 아인슈타인 방정식과의 연관성을 알아보겠습니다.

구형 우주란 \emph{공간}이 구형이라는 뜻이지, 시공간이 구형이라는 뜻은 아닙니다. 또 다른 가능성은 균일한 음의 곡률을 가진 경우입니다. 이 두 경우를 나란히 놓고 보면 일반 상대성이론의 역할이 분명해집니다.

\section{구와 균일한 기하학}

1차원 구는 단순히 원입니다. 이를 기술하는 가장 간단한 방법은 2차원 보조 공간에 매장하고 다음과 같이 쓰는 것입니다.
\begin{equation}
x^2+y^2=R^2.
\end{equation}
단위 1차원 구는 다음과 같습니다.
\begin{equation}
x^2+y^2=1.
\end{equation}

2차원 구도 같은 방식으로 기술하되, 이번에는 보조 좌표를 세 개 사용합니다.
\begin{equation}
x^2+y^2+z^2=R^2,
\end{equation}
단위 구의 경우에는 다음과 같습니다.
\begin{equation}
x^2+y^2+z^2=1.
\end{equation}

이 좌표들은 어디까지나 편의를 위한 보조 수단입니다. 관심 있는 기하학적 대상은 주변의 \(x,y,z\)공간이 아니라 표면 자체입니다. 1차원 구 위에 사는 생명체는 원 위에 살며, 안쪽이나 바깥쪽에 접근할 수 없습니다. 2차원 구에도 같은 생각을 적용할 수 있습니다. 구의 내부를 채우는 부분과 구 바깥은 보조적인 그림에만 속합니다.

같은 구성을 통해 3차원 구도 얻을 수 있습니다. 이를테면 \(w\)이라는 좌표를 하나 더 도입해 다음과 같이 씁니다.
\begin{equation}
x^2+y^2+z^2+w^2=R^2,
\end{equation}
단위 구의 경우에는 다음과 같습니다.
\begin{equation}
x^2+y^2+z^2+w^2=1.
\end{equation}
3차원 구는 시각화하기가 더 어렵습니다. 보조 공간이 4차원이 되기 때문입니다. 하지만 수학은 그리 더 어려워지지 않습니다.

표준 표기는 \(\Omega_n\)이며, \(n\)-구면에 대해 사용한다. 따라서 \(\Omega_1\)은 1-구면, \(\Omega_2\)은 2-구면, \(\Omega_3\)은 3-구면이다.

\section{구 위의 계량}

더 유용한 기술 방법은 내재적인 방식입니다. 구 자체에 놓인 좌표를 사용해 구를 기술합니다.

단위 1차원 구는 좌표 \(\phi\) 하나면 충분하며, \(0\leq \phi \leq 2\pi\)입니다. 계량은 다음과 같습니다.
\begin{equation}
ds^2=d\phi^2.
\end{equation}
이를 다음과 같이 쓰기도 합니다.
\begin{equation}
d\Omega_1^2=d\phi^2.
\end{equation}

단위 2차원 구는 1차원 구들이 이어진 형태로 기술할 수 있습니다. 각 원을 따라 도는 각도는 여전히 \(\phi\)이지만, 구의 한쪽 끝에서 다른 쪽 끝까지 이어지는 두 번째 각도 \(\theta\)도 있습니다. 두 번째 각도는 \(0\)에서 \(\pi\)까지 변합니다. \(\theta\) 방향의 운동은 \(d\theta^2\)을 기여합니다. \(\phi\) 방향으로 도는 운동은 반지름이 \(\sin\theta\)인 원 위에서 일어납니다. 따라서
\begin{equation}
d\Omega_2^2=d\theta^2+\sin^2\theta\,d\phi^2
\end{equation}
또는
\begin{equation}
d\Omega_2^2=d\theta^2+\sin^2\theta\,d\Omega_1^2.
\end{equation}

3차원 구도 같은 방식으로 얻을 수 있으며, 이번에는 2차원 구들이 이어진 형태입니다. 좌표 \(\alpha\)를 하나 더 도입하면 다음과 같습니다.
\begin{equation}
d\Omega_3^2=d\alpha^2+\sin^2\alpha\,d\Omega_2^2.
\end{equation}
이를 전개하면 다음과 같습니다.
\begin{equation}
d\Omega_3^2=d\alpha^2+\sin^2\alpha\left(d\theta^2+\sin^2\theta\,d\phi^2\right).
\end{equation}
여기서 \(\alpha\)는 크기가 0인 한쪽 끝에서 다른 쪽 끝까지 변하므로 \(0\leq \alpha \leq \pi\)입니다.\editorialnote{00:17:24에서 오디오는 잠시 ``zero to two pi,''처럼 들리지만, 주변의 구성과 명시적으로 나타난 \(\sin^2\alpha\) 인자를 고려하면 새 좌표는 더 낮은 차원의 구들로 이루어진 계열에서 크기가 0인 한쪽 끝과 다른 쪽 끝 사이를 변해야 합니다.}

일반적인 재귀 형태는 다음과 같습니다.
\begin{equation}
d\Omega_n^2=d\alpha^2+\sin^2\alpha\,d\Omega_{n-1}^2.
\end{equation}

이 관점을 받아들이면 매장 그림은 잊어도 됩니다. 1차원 구는 좌표 하나로, 2차원 구는 좌표 두 개로, 3차원 구는 좌표 세 개로 기술하는 식입니다. 기하학은 그 자체로 성립합니다.

이 기하학적 대상이 특별한 이유는 동질적이기 때문입니다. \(\sin^2\theta\)를 양 끝에서 여전히 0이 되는 다른 함수로 바꾸어 계량을 변형할 수는 있습니다. 그러면 타원체 같은 변형된 닫힌 표면을 얻게 됩니다. 하지만 더는 균일하지 않습니다. 구가 동질적인 이유는 모든 점이 다른 모든 점과 같기 때문입니다.

단위 반지름 대신 반지름이 \(A\)인 구를 원한다면, 단위 계량에 \(A^2\)을 곱합니다.
\begin{equation}
ds^2=A^2\,d\Omega_3^2.
\end{equation}
계량은 거리의 제곱을 나타내므로 제곱이 필요합니다.

0차원 구에 대해 짧게 덧붙이겠습니다. 한 단계 더 거슬러 올라가면 다음을 얻습니다.
\begin{equation}
x^2=1,
\end{equation}
해는 두 점 \(x=\pm 1\)입니다. 따라서 단위 0차원 구는 두 점뿐입니다.

\section{구형 프리드만-르메트르-로버트슨-워커 우주}

이제 관심 대상은 공간만이 아니라 시공간입니다. 여기서 다룰 그림은 반지름이 시간에 따라 변하는 3차원 구입니다.

계속 염두에 두어야 할 주의점이 하나 있습니다. 시간은 다릅니다. 팽창하는 우주를 시간이라는 좌표가 매장 좌표 하나의 역할을 더하는 고차원 구라고 생각해서는 안 됩니다.

이 점을 염두에 두면 시공간 간격은 다음과 같습니다.
\begin{equation}
d\tau^2=dt^2-a(t)^2\,d\Omega_3^2.
\end{equation}
이것이 공간이 구형인 우주의 FRW 계량입니다. 정적인 구에 비해 새로 추가된 것은 반지름이 시간에 따라 변하는 함수 \(a(t)\)가 되었다는 점뿐입니다.

우리가 한 극점에 있도록 좌표를 정하면 편리합니다. 그 점에 특별한 의미가 있는 것은 아니며, 단지 하나의 약속일 뿐입니다. 그러면 \(\alpha\)를 단위 구 위에서 우리로부터 바깥쪽으로 향하는 방사 좌표로 생각할 수 있습니다. 이는 킬로미터 단위의 방사 거리가 아닙니다. 단위 구 위의 무차원 각도 척도이지만, 기하학에서는 방사 거리 역할을 합니다.

은하나 별 같은 국소적인 덩어리를 제외하면, 이것이 공간이 구형인 FRW 우주의 기하학 전부입니다. 공간 곡률은 양수이며, 부호 관례는 다음과 같이 씁니다.
\begin{equation}
k=+1.
\end{equation}
그 자체가 ``플러스 1''이라는 뜻은 아닙니다. 공간 곡률이 양수라는 뜻일 뿐입니다.

\section{구에서 바깥을 바라보기}

구 위의 한 점에서 바깥을 바라보면 이런 구는 어떻게 보일까요? 여기서 입체 사영을 도입합니다.

무한 평면에 남극이 접하도록 2차원 구를 놓습니다. 구 위의 한 점을 고릅니다. 북극에서 그 점을 지나는 직선을 긋고, 평면과 만날 때까지 연장합니다. 이렇게 구의 각 점을 평면의 한 점에 대응시킵니다. 남극은 평면의 원점에 대응합니다. 북극은 예외적인 점입니다. 서로 다른 방향에서 북극에 다가가면 그 상은 방향에 따라 다르게 무한대로 뻗어 나갑니다. 따라서 북극은 무한대의 점으로 취급합니다.

같은 크기의 ``대륙''을 구의 서로 다른 위치에 놓으면, 북극 쪽으로 갈수록 그 상은 평면에서 더 커집니다. 따라서 남극에 있는 관측자는 먼 물체가 평평한 공간에서보다 더 큰 각도로 보인다고 관측할 것입니다. 양의 곡률에서는 알려진 크기의 물체가 유클리드 기하학에 비해 너무 크게 보입니다.

이것이 관측에서 핵심이 되는 점입니다. 곡률은 거리, 크기, 각 크기 사이의 관계로 추론합니다. 은하 자체는 아주 먼 거리에서 표준 자로 쓰기에는 적합하지 않지만, 논리는 바로 이와 같습니다.

\begin{classroomqa}
\lecturetimestamp{00:45:54}
\audiencequestion{각 크기만이 아니라 거리도 알아야 하나요?}
\lecturerresponse{그렇습니다. 각도만으로는 충분하지 않습니다. 각도와 거리가 모두 필요합니다.}
\end{classroomqa}

\begin{classroomqa}
\lecturetimestamp{00:46:18}
\audiencequestion{구 표면의 원은 사영해도 원으로 나타나나요?}
\lecturerresponse{그렇습니다. 원은 사영해도 원으로 나타납니다. 이것은 입체 사영의 특별한 성질 중 하나입니다.}
\end{classroomqa}

\begin{classroomqa}
\lecturetimestamp{00:47:02}
\audiencequestion{사영해도 대륙의 모양이 그대로 유지되나요?}
\lecturerresponse{서로 다른 두 가지 명제가 있습니다. 하나는 원이 사영해도 원이 된다는 것입니다. 다른 하나는 그보다 약한 명제로, 각도가 보존되기 때문에 충분히 작은 모양은 국소적인 형태를 유지한다는 것입니다. 이 사상은 등각 사상입니다. 상대적인 크기는 보존하지 않지만 각도는 보존합니다.}
\end{classroomqa}

\section{균일한 음의 곡률}

균일한 음의 곡률을 가진 경우는 그 계량을 통해 도입합니다.

구에 대응하는 음의 곡률 공간은 쌍곡공간입니다. 앞에서와 같은 재귀적인 방식으로 계량을 얻으며, 이때 \(\sin\alpha\)를 \(\sinh\alpha\)로 바꿉니다.
\begin{equation}
d\mathcal H_n^2=d\alpha^2+\sinh^2\alpha\,d\Omega_{n-1}^2.
\end{equation}
더 낮은 차원의 각도 부분은 단위 구 계량 \(d\Omega_{n-1}^2\)으로 그대로 남습니다.

\begin{classroomqa}
\lecturetimestamp{00:55:51}
\audiencequestion{계량에서 더 낮은 차원의 미분 항은 여전히 구형 계량인가요?}
\lecturerresponse{그렇습니다. 더 낮은 차원의 부분은 보통의 단위 구 계량입니다. 달라지는 것은 그 앞에 붙는 \(\sinh^2\alpha\) 인자입니다.}
\end{classroomqa}

쌍곡사인은 다음과 같습니다.
\begin{equation}
\sinh\alpha=\frac{e^\alpha-e^{-\alpha}}{2}.
\end{equation}
마이너스 부호가 중요합니다. 강의에서는 학생의 지적 직후 이를 바로잡습니다.

좌표의 범위에도 중요한 차이가 하나 있습니다. 구의 경우 새 좌표는 한쪽 끝에서 다른 쪽 끝까지 변하며 \(\pi\)에서 멈춥니다. 여기서는
\begin{equation}
0\leq \alpha<\infty.
\end{equation}
보통 사인 함수는 증가했다가 다시 0으로 돌아옵니다. 쌍곡사인 함수도 \(\alpha=0\) 근처에서는 같은 방식으로 시작하지만, 그 뒤에는 지수적으로 증가합니다.

따라서 바깥쪽으로 이동할수록 우리를 둘러싼 더 낮은 차원의 구들은 점점 더 커집니다. 이 기하학에서는 반지름을 늘려 가며 보아도 하늘이 먼저 커졌다가 다시 작아지지 않습니다. 계속 커집니다.

공간 차원이 3일 때 이에 대응하는 FRW 형식은 다음과 같습니다.
\begin{equation}
d\tau^2=dt^2-a(t)^2\,d\mathcal H_3^2,
\end{equation}
그리고 부호 표기는 다음과 같습니다.
\begin{equation}
k=-1.
\end{equation}

\section{원판 그림과 그 의미}

음의 곡률은 에셔의 작품, 특히 \emph{제한 원 4}으로 유명해진 원판 그림을 통해 직관적으로 이해할 수 있습니다. 입체 사영과의 유사성은 유용하지만, 그 작동 방식은 정반대입니다.

구의 경우 실제 기하학은 유한하지만 사영 그림은 무한해 보입니다. 음의 곡률에서는 실제 기하학이 무한하지만 원판 그림은 유한해 보입니다. 같은 크기의 ``대륙''이나 ``은하''는 경계에 가까워질수록 점점 작게 보이지만, 기하학적으로는 모두 같은 크기입니다. 같은 크기의 조각을 몇 개 지나왔는지 세어 거리를 잰다면, 경계까지의 거리는 무한합니다.

따라서 음의 곡률에서는 알려진 크기의 물체가 평평한 공간에 비해 너무 작게 보입니다. 물체가 차지하는 각도도 너무 작습니다.

기하학은 여전히 동질적입니다. 그림에서는 바깥쪽 조각들이 작아 보이지만, 관측자를 다른 위치로 옮겨도 내재적인 기하학은 똑같이 보일 것입니다.

\begin{classroomqa}
\lecturetimestamp{01:06:42}
\audiencequestion{그 원들 중 어디에 있는 관측자라도 비슷한 모습을 보게 되나요?}
\lecturerresponse{그렇습니다. 정확히 같은 모습을 보게 됩니다. 공간은 동질적입니다.}
\end{classroomqa}

\begin{classroomqa}
\lecturetimestamp{01:07:03}
\audiencequestion{구형인 경우에는 먼 물체가 너무 크게 보이고, 여기서는 너무 작게 보이나요?}
\lecturerresponse{그렇습니다. 양의 곡률에서는 알려진 크기의 물체가 평평한 공간에 비해 너무 큰 각도를 차지합니다. 음의 곡률에서는 너무 작은 각도를 차지합니다.}
\end{classroomqa}

\begin{classroomqa}
\lecturetimestamp{01:08:09}
\audiencequestion{천문학자가 관측으로 그 모습을 재구성할 수 있다는 점에서 입체 사영 방식의 사영이 중요한가요?}
\lecturerresponse{그렇습니다. 천문학자는 사실상 하늘의 삼각형을 측량하는 측량사입니다. 당장의 기하학적 요점에 집중하려면 시간 의존성은 잠시 제쳐 두고 각도와 거리를 보면 됩니다.}
\end{classroomqa}

국소적인 그림도 유용합니다. 음의 곡률을 가진 작은 조각은 보통의 3차원 공간에 안장 표면으로 그릴 수 있습니다. 하지만 그렇게 그릴 수 있는 것은 작은 조각뿐입니다. 쌍곡평면 전체를 3차원 유클리드 공간에 매끄러운 보통의 표면으로 전역적으로 매장할 수는 없습니다.

\section{평평한 경우}

양의 곡률과 음의 곡률을 가진 경우 사이에는 평평한 공간이 있습니다. 이를 극한으로 생각할 수도 있지만, 보통의 평평한 유클리드 공간으로 생각하는 것만으로 충분합니다.

그 FRW 형식은 다음과 같습니다.
\begin{equation}
d\tau^2=dt^2-a(t)^2\left(dx^2+dy^2+dz^2\right),
\end{equation}
그리고
\begin{equation}
k=0.
\end{equation}

이것이 세 가지 동질적인 FRW 기하학입니다.
\begin{align}
k=+1 &: \text{positive spatial curvature},\\
k=0 &: \text{flat space},\\
k=-1 &: \text{negative spatial curvature}.
\end{align}

\begin{classroomqa}
\lecturetimestamp{01:15:07}
\audiencequestion{기하학을 평면에 사영하고 신호의 경로를 따라간다면, 사영 그림에서 빛의 속도는 위치에 따라 달라지나요?}
\lecturerresponse{그렇습니다. 사영 그림에서는 달라집니다. 구의 사영에서는 바깥쪽으로 갈수록 빛이 점점 빨라지는 것처럼 보이고, 음의 곡률 원판에서는 가장자리에 가까워질수록 점점 느려지는 것처럼 보입니다. 그런 변화는 그림에 속할 뿐, 바탕이 되는 물리에는 속하지 않습니다.}
\end{classroomqa}

\begin{classroomqa}
\lecturetimestamp{01:17:05}
\audiencequestion{이런 방식으로 먼 천체를 관측해 평탄성을 알아낸 적이 있나요?}
\lecturerresponse{그렇습니다. 하지만 실제로 사용하는 것은 각도 정보입니다. 사영 그림에서 빛이 얼마나 빠르게 움직이는 것처럼 보이는지를 재서 평탄성을 알아내지는 않습니다.}
\end{classroomqa}

\section{우주론 방정식으로 돌아가기}

휴식 시간 뒤에는 앞서 살펴본 우주론 방정식과 아인슈타인 방정식의 연관성을 다룹니다. 이 시점에 칠판에는 평평한 공간 계량의 일부가 아직 남아 있습니다.

\begin{figure}[htbp]\centering
\includegraphics[width=0.72\textwidth]{lecture_05_figure_02.png}\caption{평평한 FRW 계량의 일부.}\lectureframe{01:18:18}\end{figure}

먼저 뉴턴 역학을 아주 간단히 복습하겠습니다. 우리를 중심에 둡니다. 허블 법칙에 따라 멀어지는 은하를 살펴봅시다. 구대칭 분포에 대한 뉴턴 정리에 따르면, 그 은하에 작용하는 힘에는 은하까지의 거리와 같은 반지름을 가진 구 안의 질량만 기여합니다. 그러면 가능성은 세 가지입니다. 은하가 탈출 속도보다 빠르게 움직이거나, 탈출 속도보다 느리게 움직이거나, 정확히 탈출 속도로 움직입니다.

다음으로 기억할 점은 보통의 비상대론적 물질에서는 밀도가 질량을 \(a^3\)로 나눈 것에 비례하고, 복사에서는 \(1/a^4\)에 비례하며, 진공 에너지의 밀도는 구가 변해도 변하지 않는다는 것입니다.

앞서 다룬 뉴턴 역학에서는 우주론 방정식에 밀도 항이 있었고, 정확히 탈출 속도에 해당하는 경우가 아니라면 \(1/a^2\)에 비례하는 상수항이 하나 더 있었습니다. 앞선 논의에서는 대부분 그 항을 0으로 두었습니다. 경험적으로 그 값은 작지만, 개념적으로는 제외할 이유가 없습니다.

\section{아인슈타인 방정식과 \(00\) 성분}

다음 단계는 아인슈타인 장 방정식입니다. 왼쪽에는 시공간 곡률로 이루어진 기하학이 놓이고, 오른쪽에는 중력장의 원천인 에너지-운동량 텐서가 놓입니다. 강의에서 사용한 정규화에서는 다음과 같이 씁니다.
\begin{equation}
R^{\mu\nu}-\frac{1}{2}g^{\mu\nu}R=4\pi G\,T^{\mu\nu}.
\end{equation}
\editorialnote{칠판에 보이는 식에서는 인자 \(\tfrac12\)가 읽히지 않지만, 01:33:36 부근의 구두 정정으로 그 인자를 복원할 수 있습니다. 여기 적은 방정식은 그 정정을 반영하면서 강의의 \(4\pi G\) 정규화를 유지합니다.}

서로 독립적인 열 개 성분을 모두 쓸 필요는 없습니다. 이 시점에는 이미 많은 조건,---동질성, FRW 형식, 시간에 따라 변하는 단일 척도 인자---가 가정되어 남은 임의 함수는 \(a(t)\) 하나뿐입니다. 자연스럽게 고를 성분은 시간-시간 성분입니다. \(T^{00}\)이 에너지 밀도 \(\rho\)이기 때문입니다.

\begin{figure}[htbp]\centering
\includegraphics[width=0.62\textwidth]{lecture_05_figure_05.png}\caption{아인슈타인 장 방정식과 그 \(00\) 성분.}\lectureframe{01:30:54}\end{figure}

따라서 초점은 다음 식입니다.
\begin{equation}
R^{00}-\frac{1}{2}g^{00}R=4\pi G\,\rho.
\end{equation}

오른쪽은 에너지 밀도입니다. 왼쪽은 기하학입니다.

\begin{classroomqa}
\lecturetimestamp{01:39:09}
\audiencequestion{에너지-운동량 텐서의 다른 성분들은 어떻습니까?}
\lecturerresponse{일반적으로 사라지지 않습니다. 공간-공간 성분은 압력을 나타낼 수 있습니다. 물질 지배의 경우 압력은 영이고, 복사 지배의 경우 압력은 에너지 밀도의 삼분의 일이며, 진공 에너지의 경우에는 \(P=w\rho\)를 쓰고 \(w=-1\)로 둡니다.}
\end{classroomqa}

\begin{classroomqa}
\lecturetimestamp{01:40:24}
\audiencequestion{그렇다면 물질만 있다면 \(T^{\mu\nu}\)의 다른 성분들은 사라집니까?}
\lecturerresponse{압력이 없는 물질이라면 그렇습니다. 압력 성분은 사라집니다.}
\end{classroomqa}

\section{아인슈타인 계산에서 얻는 것}

칠판에서 곡률 계산 전체를 전개하지는 않습니다. 중요한 것은 FRW 계량을 아인슈타인 텐서에 대입하면 무엇이 나타나는가입니다.

한 종류의 항은 척도 인자가 시간에 의존한다는 사실만으로 생깁니다. 공간이 평평하더라도 시간에 따라 변하는 척도 인자는 여전히 시공간 곡률을 만듭니다. 이 기여는 다음에 비례합니다.
\begin{equation}
\left(\frac{\dot a}{a}\right)^2.
\end{equation}

두 번째 종류의 항은 공간 곡률 자체에서 생깁니다. 척도 인자가 시간에 따라 변하지 않더라도 곡률을 가진 공간은 여전히 시공간 곡률에 기여합니다. 균일한 곡률을 가진 단순한 기하학에서는 곡률이 곡률 반지름의 제곱에 반비례하므로, 다음에 비례하는 항을 기여합니다.
\begin{equation}
\frac{k}{a^2},
\end{equation}
구에서는 부호가 양수이고, 쌍곡 경우에는 음수이며, 평평한 공간에서는 0입니다.

\begin{classroomqa}
\lecturetimestamp{01:36:53}
\audiencequestion{곡률 반지름 같은 인자도 있어야 하지 않나요?}
\lecturerresponse{그것이 이미 \(1/a^2\) 인자입니다. 이런 단순하고 균일한 곡률 기하학에서는 곡률이 곡률 반지름의 역제곱일 뿐입니다.}
\end{classroomqa}

계수는 계산을 진행하면서 정리하며, 앞서 얻은 뉴턴 역학의 답과 일부 대조합니다. 최종적으로 얻는 것은\footnote{이 프리드만 형식의 식은 계수에 관한 구두 설명과 자체 정정을 한데 모으고, 01:32:35--01:36:53 무렵 앞서 얻은 뉴턴 역학의 결과와 비교한 내용도 반영합니다.}
\begin{equation}
\left(\frac{\dot a}{a}\right)^2+\frac{k}{a^2}=\frac{8\pi G}{3}\rho,
\end{equation}
또는 동등하게
\begin{equation}
\left(\frac{\dot a}{a}\right)^2=\frac{8\pi G}{3}\rho-\frac{k}{a^2}.
\end{equation}

이 형식은 세 가지 FRW 기하학을 앞서 다룬 우주론적 동역학과 연결합니다.

\begin{classroomqa}
\lecturetimestamp{01:37:52}
\audiencequestion{\(k\)를 재정의하면 \(a\)는 어떤 값이든 될 수 있나요?}
\lecturerresponse{평평한 경우에는 \(a\)에 대해 유일하게 정해지는 자연스러운 척도가 없으므로 재조정할 수 있지만, \(k\)는 여전히 0입니다. 구형인 경우와 음의 곡률인 경우에는 자연스러운 곡률 척도가 있습니다. \(a(t)\)를 그 자연스러운 곡률 반지름으로 잡으면, \(k\)는 각각 \(+1\) 또는 \(-1\)일 뿐입니다.}
\end{classroomqa}

\begin{classroomqa}
\lecturetimestamp{01:41:00}
\audiencequestion{그 최종 방정식에서 \(k\)는 \(+1\), \(-1\), \(0\) 중 무엇인가요?}
\lecturerresponse{그렇습니다. 구형인 경우에는 \(+1\), 쌍곡 경우에는 \(-1\), 평평한 공간에서는 \(0\)입니다.}
\end{classroomqa}

뉴턴 역학의 설명은 우리를 잘못된 방향으로 이끌지 않았습니다. 같은 기본 방정식을 주었습니다. 일반 상대성이론이 더하는 것은 그 해석입니다. 뉴턴 역학의 관점에서 추가 상수는 탈출 속도보다 빠른 경우, 느린 경우, 그리고 경계에 해당하는 경우를 구분합니다. 아인슈타인의 관점에서 같은 부호는 공간 자체의 기하학을 알려 줍니다.

\begin{classroomqa}
\lecturetimestamp{01:42:46}
\audiencequestion{우주 상수를 포함하기 전에는 어떻게 되었나요?}
\lecturerresponse{우주 상수를 포함하기 전에는 탈출 속도에 따른 운동과 공간 기하학 사이에 밀접한 관계가 있었습니다. \(k=+1\)이면 공간은 양의 곡률을 가지며 운동은 탈출 속도보다 느리므로, 우주는 팽창했다가 다시 수축할 수 있습니다. \(k=-1\)이면 공간은 음의 곡률을 가지며 운동은 탈출 속도보다 빠르므로, 우주는 계속 팽창합니다. \(k=0\)이면 경계에 해당하는 경우입니다.}
\end{classroomqa}

따라서 같은 부호 \(k\)를 두 가지 방식으로 읽을 수 있습니다. 운동의 동역학적 분류로 볼 수도 있고 공간 곡률의 부호로 볼 수도 있습니다.